% SGA 3, Expose V, Loop 1 English reconstruction.
% Scope: Lemma 7.3 through the finiteness argument, including its closing
% seam on component-local p. 27.
% Authority: Polo--Gille Exp5-13oct24.pdf, component-local pp. 26--27
% (combined-reader pp. 288--289), re-edited headers 276--277.
% Jacob Reinhold's CC BY 4.0 Markdown is comparison/drafting material only.

\medskip
\phantomsection\label{lem:sga3-V-7.3}%
\noindent\textbf{引理 7.3.}---%
\emph{沿用定理~\hyperref[thm:sga3-V-7.1]{7.1} 的记号与假设.
每个相对于 \(S\) 闭的点 \(z\in X_0\) 都有饱和开邻域 \(U_z\),
使由 \(X_\ast\) 在 \(U_z\) 上诱导的群胚具有拟截面.}

此断言在 \(S\) 上是局部的, 因此可假设 \(S\) 仿射且为 Noether 概形,
并将\hyperref[lem:sga3-V-7.2]{前一个引理}应用于 \((\mathbf{Sch}/S)\) 中的图
\SGAThreeDiagram[0.30\textwidth]
  {figures/exp5/p276_lemma73_base_square.png}
  {应用\hyperref[lem:sga3-V-7.2]{引理~7.2} 的图.}
  {fig:sga3-V-7.3-base-square}
  {第 V 篇, 分篇第~26 页 / 合订本第~288 页 / 重版本页眉页码~276}
于是取 \(X_0\) 的一个闭子概形 \(F\), 使
\[
  d_0\bigl(d_0^{-1}(F)\cap d_1^{-1}(z)\bigr)
\]
有限且非空, 并使 \(d_1\) 在 \(d_0^{-1}(F)\) 上的限制
在 \(d_1^{-1}(z)\) 的各点处平坦.

设 \(F_1\) 和 \(F_2\) 分别是 \(F\) 在 \(d_0\)
和 \(d_0d'_0=d_0d'_1\) 下的逆像,
并用 \(\widetilde d_0,\widetilde d_1,\ldots\)
表示由 \(d_0,d_1,\ldots\) 诱导的态射. 于是有交换图
\SGAThreeDiagram[0.67\textwidth]
  {figures/exp5/p276_lemma73_F2_diagram.png}
  {\(F_2,F_1,F\) 上的诱导图.}
  {fig:sga3-V-7.3-F-diagram}
  {第 V 篇, 分篇第~26 页 / 合订本第~288 页 / 重版本页眉页码~276}
其中左边两个方块是笛卡尔图, 第一行正合
(参见 \hyperref[sec:sga3-V-1]{\S1} 中的图
\hyperref[sga3:v:lemma:1-2:equation:0-1-2]{\((0,1,2)\)}),
而 \(q,\widetilde q\) 为结构态射.

先证明 \(F_1\) 中只有有限个点位于 \(z\) 上.%
\footnote{\label{note:sga3-V-39}编者注: 已补出细节,
并明确写出了定理~\hyperref[thm:sga3-V-7.1]{7.1}
中关于 \(d_0\) 和 \(d_1\) 的固有性假设的作用.
比较\hyperref[thm:sga3-V-8.1]{定理~8.1} 的陈述及证明,
那里不作此固有性假设.}
设 \(s\) 是 \(z\) 在 \(S\) 中的像.
由于 \(F\) 在 \(S\) 上有限型,
纤维 \(\widetilde q^{-1}(s)\) 是 Noether 概形.
由于 \(\widetilde d_0\) 固有,
\[
  \widetilde d_0\bigl(\widetilde d_1^{-1}(z)\bigr)
\]
是 \(\widetilde q^{-1}(s)\) 的一个由有限个点组成的闭子概形.
因此 (参见 EGA~I, 6.2.2), 这些点在 \(\widetilde q^{-1}(s)\) 中闭;
由于 \(F\) 在 \(X_0\) 中闭, 它们在 \(X_0\) 的纤维 \(q^{-1}(s)\) 中也闭.

设 \(y\) 是这些点中的一个.
由于 \(q^{-1}(s)\) 在 \(\kappa(s)\) 上有限型,
它分别有 \(y\) 和 \(z\) 的仿射开邻域
\[
  \operatorname{Spec}B\quad\text{和}\quad\operatorname{Spec}C
\]
其中 \(B\) 和 \(C\) 是有限生成 \(\kappa(s)\)-代数.
点 \(y\) 和 \(z\) 对应于极大理想
\[
  \mathfrak p\subset B,\qquad \mathfrak q\subset C.
\]
域 \(B/\mathfrak p\) 和 \(C/\mathfrak q\)
在 \(\kappa(s)\) 上有限, 故
\[
  (B/\mathfrak p)\otimes_{\kappa(s)}(C/\mathfrak q)
\]
是有限维 \(\kappa(s)\)-代数.
它的极大理想恰好对应于 \(X_0\times_S X_0\) 中
第二投影为 \(z\) 且第一投影为 \(y\) 的点.
因此只有有限个点 \(u\in X_0\times_S X_0\),
其第二投影为 \(z\), 第一投影属于
\[
  \widetilde d_0\bigl(\widetilde d_1^{-1}(z)\bigr).
\]
最后, 由于 \(X_1\to X_0\times_S X_0\) 的纤维有限,
每个这样的点 \(u\) 在 \(X_1\) 中仅有有限多个原像点.
这就证明了断言.
% SGA3-V-V2-EOL-PRESERVATION

% Translation coverage: entire assigned component; adjoining components retain the original statement/proof seams.
% Source lines: 92; translated body lines before this audit comment: 87.
% Math preserved: 62 inline formulas, 6 display blocks, 0 equation environments; mathematical prose translated where present.
% Labels: 2/2; hyperrefs: 7/7; SGARefLinks: 0/0; inline targets: 0/0; footnotes: 1/1.
% Figure arguments preserved; all rendered captions and natural-language source locators translated.
% Unresolved TODOs: 0; compilation and visual inspection remain with the volume integrator.
